\chapter*{使用说明}\addcontentsline{toc}{chapter}{使用说明}
本手册与全书深化稿同步，先给各章原有“练习与讨论”的参考要点，再给D1--D4深化题的解答。开放题不规定唯一措辞，但要求数据、评价和结论相互一致。具体实验分数来自本工程固定版本与划分；学生修改了数据或配置后，应依据实际结果解释，不能照抄参考数值。

本轮进行了作者自检、独立计算路径的数值检查和干净内核执行，并未由另一位作者完成独立学术审阅。课堂时长、学生实际难点和不同平台表现仍需要试教记录。程序成功运行与教学任务得到有效检验是两项不同标准。

每章课堂先要求学生根据机制写下可检验判断，再运行，最后核对。结果未按预期变化时，检查是否由样本数、任务过易、尺度或评价设计造成；不为了画出理想曲线而筛选数据。答案中标明的假设与分母是评分内容的一部分。

\chapter{机器学习的基本概念与任务}
\section*{基础练习参考要点}
\begin{enumerate}[label=\arabic*.]
\item 特征矩阵为$150\times4$；物种列是当前要预测的目标。实际预测时物种未知，将它同时作为输入会泄漏答案。
\item 训练利用已有配对记录建立规则；预测只用待判断样本的特征；评价再用保留的参考类别核对。训练知道答案，是为了对以后未知的对象作判断。
\item 不会。0、1、2只是身份编码，不表示有间距的连续测量量。任务类型由目标含义决定，不由存储为整数还是字符串决定。
\item 分类训练使用已知物种，聚类训练不使用物种标签而按特征结构分组。组号可以任意重命名，也未必与物种一一对应。
\item 样本单位决定哪些行相互关联和模型将用于何种新对象。例如每行一件材料，以测量特征预测其类别；若每件材料有多次测量，要说明按材料划分以检查新材料泛化。
\end{enumerate}
\section*{深化练习参考解答}
\begin{enumerate}[label=D\arabic*.]
\item 不必相同。甲可把一次输入—输出测量作为记录，输入为已知输入电压和可用环境量，输出电压为数值目标。乙若要判断整件传感器，可把一条曲线或其事先定义的统计量作为一个对象，目标来自独立的故障检查。不能直接把故障结论或由结论编码的列当输入；同一传感器的重复记录需按部署目标分组。评分看角色是否清楚，而非要求唯一特征。
\item 不能。确定性模型对相同输入只能给出同一类别，因此至少一条错分。可改进输入信息、检查标签或采用概率描述，但单纯增加复杂度不能从完全相同的输入中恢复被省略的信息。这个例子说明预测错误不一定都是优化没有完成。
\item 第一阶段的目标由原图自身提供，属于自监督构造；第二阶段使用给定类别，属于监督分类。两阶段都可以有数值损失，是否存在损失不能用于区分是否具有人工类别标签。图像在训练、验证之间仍需按原始对象划分，不能把同一原图的不同裁剪泄漏到两侧。
\item 形状、类型、无缺失检查可能全部通过，但特征—目标配对的含义已破坏。应保留稳定行号或对象ID，用原始表逐条核对；检查类别与明确字段的关系变化也可报警，但不能代替身份核对。随机打乱标签可以作为明确标注的负对照，不能作为原始数据保存。
\end{enumerate}

\section*{实验核对与评分重点}
对应脚本为\texttt{experiments/ch01.py}，表图保存在\texttt{results/}与\texttt{figures/}。评分先检查输入与目标、划分、模型配置和指标，再检查结果解释；只有配置一致时才比较具体数值。教师版Notebook包含本章实际输出。

\chapter{数据表示、探索与准备}
\section*{基础练习参考要点}
\begin{enumerate}[label=\arabic*.]
\item 一行是一个葡萄酒化学分析样本，13列是测量指标，类别对应三个栽培品种来源类型。目标不是质量得分或价格。
\item 前者是178行1列，只取alcohol特征；后者是1行13列，取第一条记录。前提是X为完整葡萄酒输入表。
\item 均值4、总体方差$8/3$、标准差$\sqrt{8/3}$；训练变换为$(-\sqrt{3/2},0,\sqrt{3/2})$，约$(-1.2247,0,1.2247)$。新值8变为$4/\sqrt{8/3}=\sqrt6\approx2.4495$。
\item 缺失表示没有可用值，零表示一个已知数值。中位数填补只是可计算的处理规则，不能保证恢复被漏测的真值，还可能改变分布和相关关系。
\item 重新拟合测试部分会使用测试分布，得到另一套中心和尺度；训练好的模型面对的坐标含义发生改变。应只调用已训练变换器的transform。
\end{enumerate}
\section*{深化练习参考解答}
\begin{enumerate}[label=D\arabic*.]
\item 普通左连接产生4行，A被扩成3行。先检查主键唯一性，再在合并时将\texttt{validate}设为\texttt{"one\_to\_one"}；若辅助记录确实是一对多，则先定义聚合或保留观测层次。拆行可能让同一对象跨训练验证，也改变各对象权重。不能把合并成功当成统计样本数真的增加。
\item 均值3，方差$8/3$，标准差$\sqrt{8/3}$；训练值为$(-\sqrt{3/2},0,\sqrt{3/2})$，新值为$\sqrt6\approx2.4495$。合并后的均值4、方差5，改变了坐标规则并使用了新数据分布，不能称为原先冻结的变换。数值变换本身可算，但评价协议不同。
\item 缺失本身可能反映测量条件。可用训练中位数填补数值，同时增加缺失指示列；指示列为1表示原记录缺失。但模型也可能学习仪器或处理流程的代理信息，换设备或缺失机制变化时失效。应检查分组误差、缺失率变化，不能认定缺失指示永远改善。
\item 原始距离中该维平方贡献乘以$10^6$，可能改变邻居。训练均值和标准差同步缩放后，标准化坐标保持不变，忽略舍入时距离不变。精确单特征阈值树的训练候选分组由顺序决定，正线性缩放保持分组且阈值同步缩放。有限精度、并列和具体实现仍需检查；不能推广为任何非线性变换对所有新样本预测完全不变。
\end{enumerate}

\section*{实验核对与评分重点}
对应脚本为\texttt{experiments/ch02.py}，表图保存在\texttt{results/}与\texttt{figures/}。评分先检查输入与目标、划分、模型配置和指标，再检查结果解释；只有配置一致时才比较具体数值。教师版Notebook包含本章实际输出。

\chapter{模型如何学习：损失函数与优化}
\section*{基础练习参考要点}
\noindent 本章继承已复核的\texttt{chapter03\_material/teacher\_notes.txt}与教师版Notebook中的完整八题参考答案。关键评分点如下（第2题补充完整数值）：\par
\begin{enumerate}[label=\arabic*.]
\item 区分训练获得的参数 $a,b$ 与学习率、迭代次数、批量大小等超参数；改变学习率改变优化过程，增加二次项改变模型家族。
\item 残差为1、0、$-1$，平均残差为0，MSE为$2/3$，半MSE为$1/3$；梯度为$(-2/3,0)$，学习率0.3同步更新后$a=1.2,b=1$，新半MSE为$0.64/3\approx0.21333$。两个梯度均在旧参数处计算。
\item 对 $J=(w-3)^2/2$，$J'(w)=w-3$；从 $w_0=0$、$\eta=0.5$ 得 $w_1=1.5$、$w_2=2.25$。$\eta=2$ 时 $w_1=6,w_2=0$，损失不降而循环振荡。
\item 顺序更新与同时更新使用的梯度位置不同；即使最终损失都下降，也不能证明轨迹、更新规则或收敛性质相同。
\item $30$条数据遍历$80$轮，全批量$80$次、小批量（每批6）$400$次、逐样本$2400$次更新；尾部不足整批时按实际样本数求平均。
\item 仅缩小学习率可能让部分参数移动过慢；标准化改变数值条件但不增加信息。参数需换回原尺度后比较；均值方差只由训练记录估计。
\item 报告五个种子的全部运行结果和变动幅度，不得挑选最好一次代替随机性评价；指定种子保证可复现运行，不保证泛化。
\item 专业案例须列样本单位、输入、目标、仪器/测量误差、重复观测、外推限制和科学含义；低训练损失不能直接证明物理机制。
\end{enumerate}
\section*{深化练习参考解答}
\begin{enumerate}[label=D\arabic*.]
\item 梯度$(4,1)$，更新为$(0.2,0.8)$，新目标$\frac12(4\times0.04+0.64)=0.4$，原为2.5。两方向误差因子分别$1-4\eta$、$1-\eta$，要求绝对值均小于1，得到$0<\eta<0.5$。结论针对该正定二次目标，不能直接用于所有非凸函数。
\item 无正则或整体目标统一乘5时，最优点不变，梯度乘5；学习率除5可保持同一起点的普通梯度下降轨迹。若只把数据损失乘5而惩罚不变，就降低了相对正则化强度，最优参数通常改变。比较软件默认值时必须明确各项的归一化。
\item 平均损失是标量，梯度是向量，不能互称估计。若固定当前参数并从全部样本独立均匀抽取批次，有限样本平均梯度满足$\mathbb E[\nabla R_{\mathcal B}(\theta)\mid\theta]=\nabla R(\theta)$。随机重排一轮时，后续参数与已见样本有关，条件无偏性不能不加说明地逐步套用。
\item 可能已接近驻点、学习率太小导致几乎不动、或浮点精度/记录精度掩盖变化；也可能处于平坦但非理想区域。检查梯度范数、参数改变量、更多有效位、不同合理步长及简单模型的参照解。只有“图平”不足以证明全局最优或良好泛化。
\end{enumerate}

\section*{实验核对与评分重点}
对应脚本为\texttt{experiments/ch03.py}，表图保存在\texttt{results/}与\texttt{figures/}。评分先检查输入与目标、划分、模型配置和指标，再检查结果解释；只有配置一致时才比较具体数值。教师版Notebook包含本章实际输出。

第3章较完整的原生直线与批量训练材料还保留在\texttt{chapter03\_material/}，不要把其中半MSE约定与新增全MSE例子的梯度、学习率混用。

\chapter{模型评价、选择与泛化}
\section*{基础练习参考要点}
\begin{enumerate}[label=\arabic*.]
\item 原TP=3、FN=1、FP=2、TN=4。改后TP=3、FN=1、FP=1、TN=5；准确率0.8、精确率0.75、召回率0.75。
\item 不能。用途由实际使用过程决定；根据test分数多次改阈值、选特征，它就参与了方案选择而不再是独立测试。
\item 残差为0、0、2；MAE为$2/3$，RMSE为$\sqrt{4/3}\approx1.1547$，均与目标同单位。
\item 验证编号可依次为1--4、5--8、9--12，训练是其余8条。每折只能由当前训练记录学习标准化参数，否则把验证分布引入了拟合。
\item 同一天体的记录共享对象属性等信息；随机拆行可能让训练见过同一天体。预测新天体应按对象分组，评价已有天体未来记录则还要按时间定义任务。
\end{enumerate}
\section*{深化练习参考解答}
\begin{enumerate}[label=D\arabic*.]
\item 不能。需保存相同对象的参考标签、两个预测、分组与错误交叉情况；比较同一记录上的配对差异，而非把两组独立指标处理。可用预先规定的重复开发划分或适当的配对重采样评估波动，重复对象按对象重采样。不要在该验证集上无限调整后仍称其为独立证据。
\item 正类100个，TP90；负类9900个，FP495。精确率$90/(90+495)=2/13\approx0.15385$。即使大部分背景被正确排除，稀有目标候选仍可被误报主导。假设这些率能迁移到该总体；不同样本选择会改变实际比例。
\item 没有。验证折的标签参与了最初选特征。每个评估训练折内都要重新选特征；若特征数量、算法或参数还需选择，可在外层训练部分做内层验证，外层保留只评价整个选择流程。最终部署模型可在全部开发数据上重做选择，不能直接把外层预测当训练输入。
\item 两者评价目标不同。随机拆行允许训练看见同一对象的其他重复信息，接近“已有对象的另一条记录”；分组更接近“新对象”。模拟标签按对象随机给出，新对象没有可学习规律，低分正暴露了限制。应根据实际部署选择，不能在二者中挑高分而不说明对象是否见过。
\end{enumerate}

\section*{实验核对与评分重点}
对应脚本为\texttt{experiments/ch04.py}，表图保存在\texttt{results/}与\texttt{figures/}。评分先检查输入与目标、划分、模型配置和指标，再检查结果解释；只有配置一致时才比较具体数值。教师版Notebook包含本章实际输出。

\chapter{线性回归与正则化}
\section*{基础练习参考要点}
\begin{enumerate}[label=\arabic*.]
\item $\hat y=w_0+w_1x+w_2x^2$ 对参数 $w_j$ 仍线性，虽然相对于输入 $x$ 呈曲线；模型的线性性指参数进入预测的形式。
\item 标准化后系数单位改变；若使用完全对应的预测函数且求得相同最优解，可以得到相同预测，但带正则化且不调整惩罚时解可能改变。
\item 不能。相关特征可能被替代；Lasso选择还受量纲、$\alpha$、数据波动影响。特征选择不等于确定因果机制。
\item RMSE中较大残差会平方放大影响；MAE线性累计误差。删除极端值应检查测量记录、质量标志和样本选择机制，不能为了提高成绩任意删除。
\item OGLE可按观测区域/周期范围留出以测试迁移或外推；Airfoil可按工况组合留出，检查未见条件。需说明样本数量、对应真实应用和可能的分布变化。
\end{enumerate}
\section*{深化练习参考解答}
\begin{enumerate}[label=D\arabic*.]
\item $\bar x=2,\bar y=10/3$，斜率$3/2$，截距$1/3$，预测$(11/6,10/3,29/6)$，按预测减目标的残差$( -1/6,1/3,-1/6)$，SSE$1/6$。对截距的一阶条件给出残差和0。它是训练优化条件，不排除残差随输入弯曲、分组偏差或验证偏差。
\item Lasso型解为$\max(2-\lambda,0)$，依次2、1.5、1、0、0。Ridge型解为$2/(1+\lambda)$，依次2、4/3、1、2/3、1/2。绝对值在零点的折角允许一段系数被压到零；平方惩罚连续缩小。本例是单参数约定，不可直接把这里$\lambda$数值套入未统一归一化的接口。
\item 预测只依赖系数和，如$(1,2)$与$(0,3)$作用相同，数据不能唯一分解贡献。Ridge可按惩罚偏好选出一个唯一解，但没有创造区分两个相同测量来源的新证据。系数稳定、数值唯一不等于物理因果可辨识。
\item OGLE是变换输入周期，目标仍为I波段视星等，输出直接具有mag单位；不是对目标再取对数。若对一般正目标取对数并拟合，指数还原涉及预测条件中位数/均值及噪声分布的区别，不能自动视为无偏均值。需清楚区分输入变换与目标变换。
\end{enumerate}

\section*{实验核对与评分重点}
对应脚本为\texttt{experiments/ch05.py}，表图保存在\texttt{results/}与\texttt{figures/}。评分先检查输入与目标、划分、模型配置和指标，再检查结果解释；只有配置一致时才比较具体数值。教师版Notebook包含本章实际输出。

\chapter{逻辑回归与概率分类}
\section*{基础练习参考要点}
\begin{enumerate}[label=\arabic*.]
\item 逻辑回归对线性得分经链接函数输出类别概率，训练使用类别标签，属于分类模型；名称沿用统计学传统。
\item $\sigma(0)=0.5$。调整阈值是在固定预测分数上改变判定，不要求重新拟合参数，但改变分类误报/漏检比例。
\item 重视正类召回率、漏检数，以及精确率、误报代价；阈值选择应通过开发集完成，并说明业务/科学目标。
\item 同一测量过程产生的变量往往相关，违背朴素独立近似；但模型仍可在某些数据上取得有效判别表现，须用独立评价核验。
\item 输入特征连接到加权求和节点，经Sigmoid输出一个二分类概率；无隐藏层。学生应标出 $w,b$ 和输出阈值的不同作用。
\end{enumerate}
\section*{深化练习参考解答}
\begin{enumerate}[label=D\arabic*.]
\item 对$p$求导为$-y/p+(1-y)/(1-p)=(p-y)/[p(1-p)]$，而$dp/dz=p(1-p)$，相乘得到$p-y$。若$z=w^Tx+b$，则权重梯度为$(p-y)x$、偏置梯度为$p-y$；对样本平均并加正则项后即可用于梯度型优化。不是所有实现都用同一种更新器。
\item 按TP、FP、FN、TN排列依次$(4,2,0,2)$、$(3,1,1,3)$、$(2,0,2,4)$；三个准确率都为6/8。精确率依次2/3、3/4、1，召回率1、3/4、1/2。准确率受两类错误共同影响，不必单调；本例展示相同准确率掩盖不同取舍。
\item 判正代价$1-p$，判负代价$9p$；前者较小当$p>0.1$，相等时按预定规则处理。阈值依赖代价定义及概率与目标总体相符；类别比例或测量条件变化、模型未校准会破坏数值含义。不能把0.1当成所有高召回任务的通用阈值。
\item 公共因子$e^{1000}$在分子分母相消，理论概率不变。计算时减去最大得分再指数，避免溢出。三类可以得到$(0.4,0.35,0.25)$，第一类仍是最大。需要拒识时应另定义阈值和评价，不可借用二分类0.5规则。
\end{enumerate}

\section*{实验核对与评分重点}
对应脚本为\texttt{experiments/ch06.py}，表图保存在\texttt{results/}与\texttt{figures/}。评分先检查输入与目标、划分、模型配置和指标，再检查结果解释；只有配置一致时才比较具体数值。教师版Notebook包含本章实际输出。

\chapter{K近邻方法}
\section*{基础练习参考要点}
\begin{enumerate}[label=\arabic*.]
\item 应写出三条训练记录坐标、目标点坐标、欧氏距离排序和前三个类别，再说明多数投票结果；正文已给五点手算例子：查询点(2,2)的三近邻距离为1、1、$\sqrt2$，类别为0、0、1，因此预测0；题目自选坐标时应另行计算。
\item 米变成毫米使该维差值扩大$1000$倍，并可能主导欧氏距离，导致邻居排序改变；需用训练折内尺度处理。
\item KNN通常存储训练实例，预测时计算相似度；逻辑回归在\texttt{fit}时估计有限参数。统一的API不等于相同训练算法。
\item 记录最近邻距离、与训练分布的比较、特征范围及是否外推；可以设置基于开发集的距离告警阈值，并说明高维距离问题。
\end{enumerate}
\section*{深化练习参考解答}
\begin{enumerate}[label=D\arabic*.]
\item 均匀投票预测0；倒数权重10、5、0.5，类别1权重10大于类别0总5.5，因此预测1。距离为0不能直接取倒数；应检查实现规定，例如只让重合记录参与并处理它们的标签冲突，不能简单把无穷值代入后继续计算。
\item 均匀均值5；权重为1和1/2，归一化后预测$(2+4)/(1.5)=4$。只要权重非负且总和为1，就位于被使用目标的最小值2与最大值8之间。这个凸组合性质说明常规近邻回归不能沿线性趋势越过邻居目标范围外推。
\item 每维差为$N(0,2)$，平方均值2、方差8；独立维相加得到均值$2p$、方差$8p$，平方距离的变异系数为$\sqrt{2/p}$。这只适用于给定独立分布的构造，用来解释一种距离集中现象，不是所有高维数据的普遍定量定律。
\item 不证明。标准化消除量纲量级差，却不保证变量相关性和任务信息相同。固定原数据划分，在训练/验证上各生成独立噪声列，用训练折拟合尺度；比较不同噪声列数量下同一K的验证结果。若还选特征数或K，应纳入开发验证，不能按最终测试表现挑选。
\end{enumerate}

\section*{实验核对与评分重点}
对应脚本为\texttt{experiments/ch07.py}，表图保存在\texttt{results/}与\texttt{figures/}。评分先检查输入与目标、划分、模型配置和指标，再检查结果解释；只有配置一致时才比较具体数值。教师版Notebook包含本章实际输出。

\chapter{决策树}
\section*{基础练习参考要点}
\begin{enumerate}[label=\arabic*.]
\item $G=1-(5/8)^2-(3/8)^2=15/32=.46875$；$H=-(5/8)\log_2(5/8)-(3/8)\log_2(3/8)\approx.954434$比特。多数类错分率为3/8；三者度量不同，不能用相同数值互换。
\item 2.6分成4A与2A+4B，下降$16/75$、错2条；2.9分成5A与1A+4B，下降.32、错1条；3.5分成6A+1B与3B，下降$54/175\approx.308571$、错1条。两种一次错分的方案其叶内比例和权重不同，因此Gini不同。
\item 利用$\sum_i(y_i-c)^2=\sum_i(y_i-\bar y)^2+n(c-\bar y)^2$证明均值最优。父SSE58，左右各2，总4，减少54；$n\Delta V=6\times9=54$。
\item 贪心根$x_1\le2.5$，左子再按$x_2\le2.5$分裂，右侧纯B；左上混合叶一A一B，至少错1。替代根$x_2\le2.5$，左按$x_1\le2.5$、右按$x_1\le1.5$，四叶纯。后剪枝只能删原树分支，不能交换根和重建另一种完整结构。
\item 整体加权下降$(20/100)\times.1=.02$；sklearn的\texttt{min\_impurity\_decrease}使用包含节点权重比例的整体下降，不能直接拿局部.1比较。
\item 单叶、两叶、六叶代价依次：$\alpha=.15$时$9.816667,.966667,.9$，选六叶；$.3$时$9.966667,1.266667,1.8$，选两叶；12时$21.666667,24.666667,72$，选单叶。这里只比较指定三棵，实际软件还可能有其他叶数的子树。
\item 同一90/30开发划分中深度3及以上验证均29/30；训练分数可继续改善。应报告平台及样本数，不声称观察到不存在的下降；过拟合风险与本次曲线证据分开。
\item 单特征重训是另一个仅用收入的模型；条件切片固定其余七列，只考察给定条件下的函数；全样本预测散点混合了其他列变化，不能把同一收入下多个预测强行连成唯一阶梯。
\item 评分检查对象、字段与单位、逐步路由、叶的训练支持及验证错误。测量扰动可能跨阈值，重复对象可能泄漏，范围外仍落入已有叶；模型规则不自动成为物理机制。
\end{enumerate}
\section*{深化练习参考解答}
\begin{enumerate}[label=D\arabic*.]
\item 2.6的下降$16/75\approx0.213333$优于3.2的$49/300\approx0.163333$；但全部候选中2.9的下降0.32更大。结论必须说明候选范围，不能用两候选比较证明全枚举最优。
\item 正线性变换保持顺序，中点和查询点同步缩放，理想精确计算下分组及路由一致。一般严格递增变换保持训练记录的候选分组，但变换后中点对应的原坐标不必等于原中点，处在相邻训练值之间的新点可能改变分支。还要考虑浮点与并列，不可用训练分组不变推断所有预测不变。
\item 无权重比例各1/2，Gini1/2；有权重比例A为.2、B为.8，Gini$.32$。加权叶概率是按指定权重的经验类别构成，不再是原始行数比例，也不自动是目标总体的已校准概率。应在报告里说明权重目的。
\item 可能记住标签噪声、重复对象泄漏、训练覆盖过窄或选择机制代理。固定验证对象，在开发训练样本上重采样，重训相同配置，比较根特征、阈值和同一验证对象预测的变化；同时区分结构变化与预测变化。仍不能由结构稳定直接推出因果机制。
\end{enumerate}

\section*{实验核对与评分重点}
对应脚本为\texttt{experiments/ch08.py}，表图保存在\texttt{results/}与\texttt{figures/}。评分先检查输入与目标、划分、模型配置和指标，再检查结果解释；只有配置一致时才比较具体数值。教师版Notebook包含本章实际输出。

\chapter{随机森林与 Bagging}
\section*{基础练习参考要点}
\begin{enumerate}[label=\arabic*.]
\item Bootstrap按有放回抽样，因此一个样本可重复、也有样本没被抽到；不同树的抽样结果独立产生但可能有共同记录。
\item Bagging通过重采样组合学习器；随机森林的树分裂通常还采用随机候选特征，从两个来源增加树之间的差异。
\item 不能。袋外只保证具体行未被该树采到；同一个天体的另一条观测可能仍在该树训练中。需要按实体分组评价。
\item 重要性描述模型在当前数据中的预测依赖或分裂行为，可能受相关变量、偏差和选择机制影响，不等于控制变量实验或因果效应。
\end{enumerate}
\section*{深化练习参考解答}
\begin{enumerate}[label=D\arabic*.]
\item 展开方差的对角与交叉项：$B^{-2}[B\sigma^2+B(B-1)\rho\sigma^2]=\sigma^2[\rho+(1-\rho)/B]$。代入为.28。模型无限多时下限为.2而非0。假设简化且协方差矩阵需有效，实际森林不具有完全相同的方差与相关性；公式说明不同误差相关带来的收益上限。
\item 每次未抽中概率$1-1/n$，连续$n$次为$(1-1/n)^n\to e^{-1}$。这是某树相对于一条记录的袋外概率，不说明不同树袋外集合独立，也不排除同一物理对象的其他记录进入训练。
\item 硬投票A两票获胜，概率平均A为$1.03/3\approx.343333$，二分类时选B。类别数字只是编码，平均数字没有所需类别含义。sklearn分类森林按平均类别概率决定类别，解释程序时应按实际聚合规则。
\item 在开发评价中比较单列置换、相关特征组联合置换、删除该列后重新训练，以及适当的条件分析；保持验证记录与流程一致。结果可以揭示当前模型和输入集合的预测依赖，但相关结构、测量与选择机制仍可能造成替代或混杂，不构成因果证明。
\end{enumerate}

\section*{实验核对与评分重点}
对应脚本为\texttt{experiments/ch09.py}，表图保存在\texttt{results/}与\texttt{figures/}。评分先检查输入与目标、划分、模型配置和指标，再检查结果解释；只有配置一致时才比较具体数值。教师版Notebook包含本章实际输出。

\chapter{梯度提升树及其应用}
\section*{基础练习参考要点}
\begin{enumerate}[label=\arabic*.]
\item 平方误差的负梯度与残差成正比；交叉熵应基于当前预测得分和概率计算对应梯度，不能无条件沿用原始 $y-\hat y$ 解释。
\item 不能。训练损失可能下降但验证误差恶化；学习率与迭代次数须联合通过开发数据评价。
\item 随机森林的树可相对独立地训练，再组合预测；提升树顺序训练，下一轮依赖前面模型的剩余误差或目标梯度。
\item 重复配方可能在训练与测试中共享材料特征，产生偏乐观评价；若目标是新配方泛化，应按配方或化学组成分组。
\end{enumerate}
\section*{深化练习参考解答}
\begin{enumerate}[label=D\arabic*.]
\item 第一棵树在前三点输出$-.75$，第四点输出2.25；新预测为$(1.375,1.375,1.375,2.875)$，MSE$.921875$。新树输出是加到已有预测上的修正，最终模型是初值与全部缩放修正的和。
\item $p=.5$，得分负梯度$y-p=.5$；类别减得分得到1，忽略了得分到概率的映射。一般损失的提升必须对模型当前所用的预测坐标求导；“残差拟合”只在适当平方损失表述下直接成立。
\item 不意味着。后续每棵树依赖先前更新后的残差或梯度，改变步长会改变之后的训练目标和划分。应固定数据、树复杂度、验证方案，分别记录逐轮训练/验证曲线与最终预测，不能只比较学习率乘树数。
\item 不能。停止轮数已由验证结果选择，最低值具有选择作用。固定选择规则与配置后，按约定在开发数据重训，再评价未参与选择的保留集；要评价整个轮数选择过程可用外层验证。运行最大轮数和early stopping的职责不同。
\end{enumerate}

\section*{实验核对与评分重点}
对应脚本为\texttt{experiments/ch10.py}，表图保存在\texttt{results/}与\texttt{figures/}。评分先检查输入与目标、划分、模型配置和指标，再检查结果解释；只有配置一致时才比较具体数值。教师版Notebook包含本章实际输出。

\chapter{支持向量机与核方法}
\section*{基础练习参考要点}
\begin{enumerate}[label=\arabic*.]
\item 支持向量是限制间隔或影响决策边界的训练样本；在软间隔条件下也可能包括间隔内部或误分类样本。
\item RBF核依赖样本间距离，不同单位和数值尺度会影响距离及核值，因此应在训练折内标准化。
\item 增大$C$通常更重视训练违例，增大$\gamma$使RBF响应更局部；二者都可能提高训练拟合但降低泛化，需验证。
\item 常见核SVM需要处理随样本数增大的核关系和优化成本；训练集极大时应与线性方法或近似核比较资源需求。
\end{enumerate}
\section*{深化练习参考解答}
\begin{enumerate}[label=D\arabic*.]
\item 对称性给$b=0$，约束$y(wx+b)\ge1$要求$w\ge1$；最小化$w^2/2$得$w=1$。支持点为$-1,1$，边界$wx+b=\pm1$间距2，单侧到分类面距离1。把w,b共同缩放不改变分类面，却改变函数得分，必须规范化后比较几何间隔。
\item $\max(0,1-yf)$分别为0、.6、1.3。得分.4为正，类别正确，但位于间隔内仍有损失。SVM不仅关心符号正确，还对间隔违例施加代价。
\item 展开为$x_1^2z_1^2+2x_1x_2z_1z_2+x_2^2z_2^2=(x_1z_1+x_2z_2)^2$。内积原值11，平方121；映射内积$1\times9+(2\sqrt2)(12\sqrt2)+4\times16=9+48+64=121$。这是核技巧的小规模示范，不说明任何相似函数都合法。
\item $\exp(-\gamma d^2)$随$\gamma$增大下降，使影响更局部。训练点可能被局部响应很好区分，但新点与各训练点核值都很小，模型难以连接邻近结构；仍受C和数据影响，并非必然结果。要联合验证并检查标准化。
\end{enumerate}

\section*{实验核对与评分重点}
对应脚本为\texttt{experiments/ch11.py}，表图保存在\texttt{results/}与\texttt{figures/}。评分先检查输入与目标、划分、模型配置和指标，再检查结果解释；只有配置一致时才比较具体数值。教师版Notebook包含本章实际输出。

\chapter{多层感知机与神经网络入门}
\section*{基础练习参考要点}
\begin{enumerate}[label=\arabic*.]
\item 在输入与输出之间增加学习到的仿射变换和非线性激活：$h=\phi(W_1x+b_1)$，输出在新表示$h$上完成；比直接逻辑回归可形成更复杂边界。
\item 多个仿射层不加非线性时可合并成一个仿射层，因此表达能力仍为线性/仿射关系。
\item 先计算$z_1=w_1x+b_1$、$h=\phi(z_1)$、$\hat y=w_2h+b_2$，再从损失对输出的导数经$w_2,\phi'(z_1)$传递到第一层，依据链式法则求梯度。
\item 反向传播计算梯度，梯度下降等优化器依据梯度更新参数；并非sklearn所有\texttt{fit()}都使用反向传播或梯度下降。
\item 不能。\texttt{loss\_curve\_}通常记录训练目标；泛化需用未参与当前训练的验证数据或独立测试评价。
\item 固定训练/验证划分、标准化、优化器和学习率、随机种子或做多种子重复；报告训练/验证指标、运行代价、收敛警告及误分类矩阵。
\item 样本量少、特征维数较高、损失表面非凸、初始化、批次顺序和正则化配置均可导致波动；需重复运行并记录分布。
\end{enumerate}
\section*{深化练习参考解答}
\begin{enumerate}[label=D\arabic*.]
\item 单隐藏层为$13\times8+8+8\times3+3=139$。两隐藏层13→8→4→3为$104+8+32+4+12+3=163$。参数数量影响表达和训练，但不是跨架构复杂度或泛化能力的唯一指标。
\item 对00、01、10、11输出0、1、1、0，给出XOR。它证明一个特定非线性表示能表达该关系，不证明从随机初值训练一定找到这些参数，也不证明更多层一定泛化更好。表达能力与优化成功是两回事。
\item 在$x=2,y=1,(w_1,b_1,w_2,b_2)=(.5,0,3,0)$得到梯度$(12,6,2,2)$；中央差分步长$10^{-6}$应接近该值。过小步长受浮点相消影响，过大步长有截断误差；ReLU折点不可微，跨折点的差分不是某个唯一导数。本例隐藏预激活为1，避开了折点。
\item 应同时报告实际错分数、各种子准确率、交叉熵、迭代次数与收敛状态；准确率只看最大类别，交叉熵还对概率分配敏感。36条验证记录较少，差异受划分和初始化影响，不能据单次微小概率损失改善宣称显著普遍优势。最终测试不得参与宽度和种子筛选。
\end{enumerate}

\section*{实验核对与评分重点}
对应脚本为\texttt{experiments/ch12.py}，表图保存在\texttt{results/}与\texttt{figures/}。评分先检查输入与目标、划分、模型配置和指标，再检查结果解释；只有配置一致时才比较具体数值。教师版Notebook包含本章实际输出。

\chapter{主成分分析与特征表示}
\section*{基础练习参考要点}
\begin{enumerate}[label=\arabic*.]
\item PCA在给定维数下选择尽量保留整体样本方差的正交方向，第一主成分对训练数据的方差解释最大。
\item 不保证。高方差可能来自无关或噪声方向，低方差特征也可能对区分类别很关键。
\item 全数据拟合PCA已借助验证折分布，可能导致评估乐观；监督预测中应置入每个训练折的\texttt{Pipeline}。
\item 若未标准化，大尺度变量可能主导方差；标准化后各列尺度接近，主成分方向可能显著变化。
\item 不可以。主成分是输入的线性组合，其关联不保证唯一物理机制；需其他观测或实验来支持解释。
\end{enumerate}
\section*{深化练习参考解答}
\begin{enumerate}[label=D\arabic*.]
\item 均值零，协方差$\frac13\left(\begin{smallmatrix}10&8\\8&10\end{smallmatrix}\right)$；特征值6与2/3，第一方向$(1,1)/\sqrt2$，第二方向$(1,-1)/\sqrt2$。丢失SSE为$(n-1)\lambda_2=2$，每样本平方距离均值.5，每坐标MSE.25。分母必须说明。
\item 单轴可比较方向的绝对内积或投影矩阵$vv^T$，整体翻号同时改变得分不影响重构。近重根时对应子空间中的正交旋转也可能产生很不同的轴；应比较子空间投影、重构及预测作用，而非把某一轴符号或细小系数差当成物理变化。
\item PCA优化输入方差，不看类别。若未标准化且第一维主导方差，保留第一主成分可能忽略低方差类别信号。可固定训练验证划分，比较原特征、PCA1和PCA2，所有尺度/PCA在训练折内拟合；标准化会改变问题但不保证最优。不要直接把高解释方差解释为高分类信息。
\item 中心化使各行之和为零，秩最多$\min(100,29)=29$。此外特征相关可进一步降低秩。输入列数是表示长度，不等于样本支持的独立方向数；小样本PCA的方向还可能不稳定。
\end{enumerate}

\section*{实验核对与评分重点}
对应脚本为\texttt{experiments/ch13.py}，表图保存在\texttt{results/}与\texttt{figures/}。评分先检查输入与目标、划分、模型配置和指标，再检查结果解释；只有配置一致时才比较具体数值。教师版Notebook包含本章实际输出。

\chapter{聚类分析}
\section*{基础练习参考要点}
\begin{enumerate}[label=\arabic*.]
\item 增大$K$提供更灵活的划分，簇内平方和通常不增；甚至每点一簇可使训练目标趋零，不能据此证明自然簇数。
\item DBSCAN的噪声标签反映局部密度与参数选择，不等于测量错误，也不必然代表稀有科学对象。
\item 不够。稳定性是内部证据；还需要视差、自行、空间位置、颜色—星等关系等与成员身份相符的外部信息。
\item 距离类方法的簇分配可能改变，尤其未处理不同特征尺度时；应说明单位、缩放方法和物理意义。
\item 无监督探索可以使用待描述的全部对象进行展示；但如果降维结果会进入某个监督模型的验证，PCA须只在相应训练折内拟合。
\end{enumerate}
\section*{深化练习参考解答}
\begin{enumerate}[label=D\arabic*.]
\item 第一次分配得到{0,1}与{4,5,9,10}，更新中心.5和7，之后不变；SSE为.5+26=26.5。另一分组{0,1,4,5}和{9,10}中心2.5和9.5，SSE17+.5=17.5更低。稳定只说明交替更新不能进一步改善当前分配，不排除其他局部解。
\item 不矛盾。轮廓系数依据簇内紧凑与簇间分离，ARI比较与参考划分的样本配对一致性。双月中K均值轮廓约.489但ARI约.257；DBSCAN轮廓约.326而ARI1。指标的偏好与非凸几何不同，不能把内部评分最大当作发现真实结构的唯一证据。
\item 不能仅靠一个非核心边界点作为核心链中继合并两簇；密度可达的扩展依赖核心点。该点可能被先到达的簇吸收，具体边界归属可随遍历顺序变化。核心结构与边界分配应区分，不把每个边界标签视为唯一科学身份。
\item 重复使用同一特征只是在检查算法已选择的性质，不是独立证据。可检查天空集中、颜色—星等位置、独立径向速度或已有成员信息。仍要考虑误差、选择函数、空间相关、外部资料质量和是否在调参时已反复使用这些信息。
\end{enumerate}

\section*{实验核对与评分重点}
对应脚本为\texttt{experiments/ch14.py}，表图保存在\texttt{results/}与\texttt{figures/}。评分先检查输入与目标、划分、模型配置和指标，再检查结果解释；只有配置一致时才比较具体数值。教师版Notebook包含本章实际输出。

\chapter{异常检测}
\section*{基础练习参考要点}
\begin{enumerate}[label=\arabic*.]
\item 不可以。异常分数表示偏离训练样本所定义的常见结构，可能由质量问题、噪声或分布差异造成；需独立证据。
\item 不意味着。5\%是人为候选预算或阈值设定，真实异常比例需要外部标注或科学调查。
\item 整体距离离群看与总体中心/邻居的偏离，局部密度异常看某个点比其邻域更稀疏；多模态数据上二者可能不同。
\item 若不区分，学生可能把已知人工注入当作真实发现，也无法明确控制实验的检出率与原始数据的科学意义。
\end{enumerate}
\section*{深化练习参考解答}
\begin{enumerate}[label=D\arabic*.]
\item 不能。经验分位数只规定有限参考样本中的排序位置，未来误报有抽样波动；分布变化、依赖样本与同数据拟合评分都影响结果。应在独立正常验证记录上估计误报，并考虑置信与部署差异。不能把contamination或quantile参数等同于真实异常比例。
\item 不能。这里人工异常被放在远离正常分布的位置，是容易的外部异常。应按预定强度、噪声、局部密度或接近边界程度设计注入阶梯，保留难例，比较误报和漏检；同时说明注入机制与真实异常物理是否相关。完美恢复一种注入不等于所有真实异常可检出。
\item 点10的2距离为9，邻居是2和1；其可达距离8和9，局部可达密度$2/17$。点1的两个邻居0、2的2距离都为2，故lrd1为1/2；点2的两个邻居1、0对应可达距离1、2，lrd2为2/3。因此LOF10=$[(1/2+2/3)/2]/(2/17)=119/24\approx4.95833$。局部密度比把对象相对于邻域自身的尺度比较，而非仅看绝对最近距离。
\item 说明评分方向、候选排序、并列处理、检查预算以及检查后的真实证据。可报告precision@20或经核查的候选数，但未知未检查对象的标签时不能给出全体召回率。前20是资源规则，不是自然界异常的定义。
\end{enumerate}

\section*{实验核对与评分重点}
对应脚本为\texttt{experiments/ch15.py}，表图保存在\texttt{results/}与\texttt{figures/}。评分先检查输入与目标、划分、模型配置和指标，再检查结果解释；只有配置一致时才比较具体数值。教师版Notebook包含本章实际输出。

\chapter{模型解释、误差分析与科学结论}
\section*{基础练习参考要点}
\begin{enumerate}[label=\arabic*.]
\item 高总体准确率可由多数类主导，稀有类召回率低意味着关键对象被遗漏；应报告分类别指标和错误例子。
\item 不能。相关特征可互相替代，置换测试可能受数据分布和特征耦合影响；预测贡献不等于物理作用或因果关系。
\item 折间标准差描述不同划分下评价指标的变动；测量误差是数据采集不确定性；预测区间则与给定预测方法及假设有关，不能等同。
\item 标注为训练支持范围外的外推，给出输入范围、质量警告和未验证条件，不以点预测数字营造确定性。
\item 表格应写明数据版本、拆分、样本数、指标定义、单位与不确定性；分组误差图应写清分组区间、每组数量和误差估计方法。
\end{enumerate}
\section*{深化练习参考解答}
\begin{enumerate}[label=D\arabic*.]
\item 不包括。它通常只描述在固定模型及当前预测记录的经验总体下，评价指标随重采样的条件波动。训练和模型选择不确定性需重训/嵌套设计，仪器变化需要外部数据。重复对象需按对象采样；指标区间也不是每个对象的预测区间。
\item 不能。两列可能互相替代，也可能逐列打乱产生不真实组合。可同一个置换索引同时打乱相关组，保持组内关系而破坏组与目标的对应；另可删去一列/一组后在开发训练中重拟合。前者解释固定模型的依赖，后者解释可重新适应时的增量信息，都不是因果作用。
\item 可能是测量噪声更大、训练支持不足、低信噪比同时对应不同目标/对象分布、目标标签误差或系统校准偏差。需比较样本数、目标范围、残差均值和离群分布，控制相关条件并核对质量。单张图不能把所有误差归因于算法，也不能证明信噪比本身的因果效应。
\item 取$k=\lceil(19+1)(1-.1)\rceil=18$，半宽$r_{(18)}$。拟合模型独立于校准标签，校准与未来样本可交换，并使用正确有限样本分位数规则。覆盖是边际意义，不保证每个子群都90\%，分布变化或重复相关破坏条件时不能照搬。
\end{enumerate}

\section*{实验核对与评分重点}
对应脚本为\texttt{experiments/ch16.py}，表图保存在\texttt{results/}与\texttt{figures/}。评分先检查输入与目标、划分、模型配置和指标，再检查结果解释；只有配置一致时才比较具体数值。教师版Notebook包含本章实际输出。

\chapter{天文与物理综合项目}
\section*{基础练习参考要点}
\begin{enumerate}[label=\arabic*.]
\item 写明X为实际预测时可用的测光字段，y来自独立可靠的光谱分类，按唯一天体ID或观测批次分组；分类评价包含混淆矩阵、各类召回及实际使用目标。
\item 明确随机划分评估的是当前混合记录分布下的表现，不能宣称可以推广到未见实验工况；应提出元数据补充和后续工况留出测试。
\item 单次最高分可能来自随机拆分、调参机会或泄漏，不能替代可复现性、误差分析、数据质量与科学解释。
\item 在开发集交叉验证中完成搜索，冻结预处理和模型后对保留集一次性评估；不能根据测试分数再重新选择再报告该分数。
\item 必须包含数据出处及版本、下载/处理代码、环境依赖、随机种子、划分策略、模型配置、运行入口、结果输出及图表重建步骤。
\end{enumerate}
\section*{深化练习参考解答}
\begin{enumerate}[label=D\arabic*.]
\item 原始为公开cepF表前90条固定节选，I目标缺失的3条排除，正周期及I可用者87条；固定种子42留20\%得到69/18。当前模型只用周期预测I，未用V，V缺失不应无理由删行。若后来引入颜色，必须重新定义可用样本与选择影响，不能悄悄改变分母。
\item 不能。可写：在指定公开目录前90条中的87条可用基频记录上，经开发集五折选择对数周期模型后，对固定18条保留记录得到约.149mag RMSE；结论仅描述该节选和划分，未系统处理消光、距离深度和样本选择。它不是全巡天代表精度或新的绝对光度标定。
\item 同批记录共享温度和零点偏移；随机拆行可能让模型见过测试批次条件，按批次留出更接近未见批次。真函数可用于教师检查生成与误差结构，但学生预测输入不能包括隐藏真值、零点偏移或事后目标。模拟真函数正确也不证明真实仪器只有这些效应。
\item 例如，第17章天文项目的结论对应\path{results/project_ogle_summary.json}；生成代码为\path{experiments/project_ogle.py}；输入为\path{data/ogle_lmc_f_excerpt.csv}；划分记录为\path{data/ogle_excerpt_split_v3.csv}。开发比较见\path{results/project_ogle_models.csv}，按其均值选模型后才评价18条保留记录。最常见混淆是把三候选中最低CV均值当独立测试，或根据18条结果再次挑模型却仍报告原测试。
\end{enumerate}

\section*{实验核对与评分重点}
对应脚本为\texttt{experiments/ch17.py}，表图保存在\texttt{results/}与\texttt{figures/}。评分先检查输入与目标、划分、模型配置和指标，再检查结果解释；只有配置一致时才比较具体数值。教师版Notebook包含本章实际输出。

真实天文节选另用\texttt{experiments/project\_ogle.py}；物理批次标定为人工模拟，二者的证据地位不同。未执行的外部全目录项目不能填入虚构结果。

\backmatter
\chapter*{试教与审阅记录建议}\addcontentsline{toc}{chapter}{试教与审阅记录建议}
课堂试教应记录学生卡住的概念、手算耗时、实验计算量、结果解释错误和环境问题。不要仅按章节字数决定删改：如果学生能运行却不能说明阈值怎样确定，就需要加强机制与算例；如果已经理解机制却被大量重复代码阻碍，应缩短辅助代码而保留分析。

技术审阅应把公式、数值例子、图、代码和参考答案作为同一组材料核对。特别关注训练目标的归一化、残差符号、样本权重、不同层级的验证用途以及对象身份。对来源数据更新或软件版本变化，保留旧版结果与差异说明，而不把历史表格静默覆盖。
